\documentclass[10pt,a4paper]{amsart}
\usepackage{amsmath,amssymb,amsthm,graphicx}
% For commutative diagrams we use Xy-pic
\usepackage[all]{xy}

% We use 2cell for 2-commutative diagrams.
\xyoption{2cell}
% Theorem environments.
%
\theoremstyle{plain}
\newtheorem{theorem}[subsection]{Theorem}
\newtheorem{proposition}[subsection]{Proposition}
\newtheorem{lemma}[subsection]{Lemma}

\theoremstyle{definition}
\newtheorem{definition}[subsection]{Definition}
\newtheorem{example}[subsection]{Example}
\newtheorem{exercise}[subsection]{Exercise}
\newtheorem{situation}[subsection]{Situation}

\theoremstyle{remark}
\newtheorem{remark}[subsection]{Remark}
\newtheorem{remarks}[subsection]{Remarks}

\numberwithin{equation}{subsection}

% Macros
%
\def\lim{\mathop{\mathrm{lim}}\nolimits}
\def\colim{\mathop{\mathrm{colim}}\nolimits}
\def\Spec{\mathop{\mathrm{Spec}}}
\def\Hom{\mathop{\mathrm{Hom}}\nolimits}
\def\Ext{\mathop{\mathrm{Ext}}\nolimits}
\def\SheafHom{\mathop{\mathcal{H}\!\mathit{om}}\nolimits}
\def\SheafExt{\mathop{\mathcal{E}\!\mathit{xt}}\nolimits}
\def\Sch{\mathit{Sch}}
\def\Mor{\mathop{\mathrm{Mor}}\nolimits}
\def\Ob{\mathop{\mathrm{Ob}}\nolimits}
\def\Sh{\mathop{\mathit{Sh}}\nolimits}
\def\NL{\mathop{N\!L}\nolimits}
\def\CH{\mathop{\mathrm{CH}}\nolimits}
\def\proetale{{pro\text{-}\acute{e}tale}}
\def\etale{{\acute{e}tale}}
\def\QCoh{\mathit{QCoh}}
\def\Ker{\mathop{\mathrm{Ker}}}
\def\Im{\mathop{\mathrm{Im}}}
\def\Coker{\mathop{\mathrm{Coker}}}
\def\Coim{\mathop{\mathrm{Coim}}}

% Boxtimes
%
\DeclareMathSymbol{\boxtimes}{\mathbin}{AMSa}{"02}

%
% Macros for moduli stacks/spaces
%
\def\QCohstack{\mathcal{QC}\!\mathit{oh}}
\def\Cohstack{\mathcal{C}\!\mathit{oh}}
\def\Spacesstack{\mathcal{S}\!\mathit{paces}}
\def\Quotfunctor{\mathrm{Quot}}
\def\Hilbfunctor{\mathrm{Hilb}}
\def\Curvesstack{\mathcal{C}\!\mathit{urves}}
\def\Polarizedstack{\mathcal{P}\!\mathit{olarized}}
\def\Complexesstack{\mathcal{C}\!\mathit{omplexes}}
% \Pic is the operator that assigns to X its picard group, usage \Pic(X)
% \Picardstack_{X/B} denotes the Picard stack of X over B
% \Picardfunctor_{X/B} denotes the Picard functor of X over B
\def\Pic{\mathop{\mathrm{Pic}}\nolimits}
\def\Picardstack{\mathcal{P}\!\mathit{ic}}
\def\Picardfunctor{\mathrm{Pic}}
\def\Deformationcategory{\mathcal{D}\!\mathit{ef}}
\usepackage{cleveref}
\crefname{originallemma}{original lemma}{original lemmas}
\Crefname{originallemma}{Original lemma}{Original lemmas}
\makeatletter
\newlabel{sites-modules-lemma-localize-morphism-ringed-sites}{{04J0}{source}{}{sourcepointer-1}{}}
\newlabel{sites-modules-lemma-localize-morphism-ringed-sites@cref}{{[originallemma][0][]04J0}{[1][0][]source}{}{}{}}
\makeatother
\title{Extension by zero along an etale morphism commutes with base change}
\author{The Stacks Project Authors}
\date{stacks-project:etale-cohomology.tex:local-snapshot-20261011; GNU FDL 1.2 or later}
\begin{document}
\maketitle
\noindent\textit{Editorial note.} All delivered mathematical text is the author's own, copied from the pinned Stacks Project chapter file etale-cohomology.tex; no statement or proof step was rewritten, reordered or removed.\par\medskip
\noindent\textit{Editorial note.} The proof is the author's complete argument: both adjunctions and the identification of the right adjoints are displayed, and the conclusion follows by uniqueness of adjoints. The Xy-pic square of the counted statement is delivered as written.\par\medskip
\noindent\textit{Editorial note.} The label lemma-shriek-base-change occurs in no other staged Stacks chapter, so no cross-chapter restatement had to be weighed for this card.\par\medskip
\noindent\textit{Editorial note.} The single labelled reference occurring inside the counted spans points to a chapter of the Project not staged in this card's source tree; it is carried in assembly.external\_references as a non-load-bearing navigation pointer recorded at the invoking line.\par\medskip
\noindent\textit{Editorial note.} The preamble is the chapter's own native preamble excerpt already used for this source: the Xy-pic setup (native preamble lines 14-18) and the native theorem declarations with the mathematical macros (native preamble lines 154-217).\par\medskip
\noindent\textit{Allowed standard prerequisites (editorial).}\begin{itemize}
\item Extension by zero j\_! along an etale morphism j : V -> X. Same chapter, the functor j\_! on abelian sheaves and on sheaves of modules, left adjoint to the restriction j\textasciicircum{}\{-1\}. Supplies the functors appearing on both sides of the asserted identity.
\item Adjunction identifying j\_! as a left adjoint. Same chapter, the adjunction between j\_! and j\textasciicircum{}\{-1\}, and between f\textasciicircum{}\{-1\} and f\_*. Supplies the argument that two functors with the same right adjoint are equal.
\item Commutation of pushforward with etale localization. Same chapter, the identification f'\_*(j')\textasciicircum{}\{-1\} = j\textasciicircum{}\{-1\}f\_* coming from the construction of pullback and pushforward. Supplies the equality of right adjoints which forces the equality of the left adjoints.
\end{itemize}\medskip
\ifdefined\hypertarget\hypertarget{sourcepointer-1}{}\fi
\noindent\textit{Source reference.} Original lemma 04J0: sources/stacks-etale-cohomology/pinned-original-source.tex, source lines 11512--11512. Cross-chapter navigation pointer invoked at proof lines 11510-11512, where the author notes that the statement holds in the general setting of a morphism of sites. The target chapter is not staged in this card's source tree, so the pointer records the invoking line; it only resolves the reference for the reader, and no bibliography citation is involved.\par\medskip
\begin{lemma}[Extension by zero commutes with base change]
\label{lemma-shriek-base-change}
Let $f: Y \to X$ be a morphism of schemes. Let $j: V \to X$ be an \'etale
morphism. Consider the fibre product
$$
\xymatrix{
V' = Y \times_X V \ar[d]_{f'} \ar[r]_-{j'} & Y \ar[d]^f \\
V \ar[r]^j & X
}
$$
Then we have $j'_! f'^{-1} = f^{-1} j_!$ on abelian sheaves and on
sheaves of modules.
\end{lemma}

\begin{proof}
This is true because $j'_! f'^{-1}$ is left adjoint to
$f'_* (j')^{-1}$ and $f^{-1} j_!$ is left adjoint to $j^{-1}f_*$.
Further $f'_* (j')^{-1} = j^{-1}f_*$ because $f_*$ commutes with
\'etale localization (by construction). In fact, the lemma holds very generally
in the setting of a morphism of sites, see
Modules on Sites, Lemma
\ref{sites-modules-lemma-localize-morphism-ringed-sites}.
\end{proof}

\end{document}
